% Volume IX: Institutional Design
% Part XVII. Designing for Delayed Truth
% Chapter 101: Procedural Regularization
% Spine: proof-spines/ch101-spine.md

\chapter{Procedural Regularization}
\label{ch:ch101}

The adapter analogy now becomes an institutional design principle. High-impact decisions should be more strongly regularized: they deserve independent review, slower update schedules, and explicit rollback mechanisms because the cost of premature or badly grounded revision is high. Low-impact and reversible decisions can tolerate faster updating because their error costs are easier to unwind.

Procedural regularization thus scales with stakes and reversibility. The point is not to make every institutional process equally slow, but to ensure that consequential judgments cannot be reconfigured by narrow evidence, transient pressure, or local enthusiasm without passing through thicker constraints. Good design decides in advance where the system should resist drift and where it may safely adapt quickly.
